\chapter{METHODOLOGY}

\section{Introduction}

The following chapter explains the research methodology and technical design choices that were made during TerraNode's development. It encompasses research design approach, architecture design of the platform, the data modelling strategy, the smart contract logic and the predictive analytics engine. In each section, the reasons for the selection of the chosen approach is included, as well as, where applicable, the mathematical representations that lie behind the mathematical processes in the system.

\section{Research Design}

The research approach in this study is Design Science Research (DSR), suitable for an information technology project to design and test a new artefact for addressing an existing problem in an organization. It is important to note that in the DSR paradigm the research process is divided into two complementary activities: artefact construction and artefact evaluation of the artefact in comparison to the requirements derived from the problem space.

The research was carried out in four stages:

\begin{enumerate}[label=(\roman*)]
    \item \textbf{Problem Identification and Requirements Analysis:} The problems and requirements of the current agricultural data management systems were identified in Chapter Two as part of the literature review. The research questions and objectives in Chapter One were used to translate into functional and non-functional requirements.
    \item \textbf{System Design:} Architecture, data models, API specifications and smart contract structures have been designed to meet the identified requirements.
    \item \textbf{Implementation:} It was implemented as a working prototype using technology stack outlined in Section~\ref{sec:tech_stack}.
    \item \textbf{Evaluation:} The system developed was tested and evaluated with the defined requirements and result has been given in Chapter~Four.
\end{enumerate}

\section{System Requirements}

\subsection{Functional Requirements}

The functional requirements specify what TerraNode needs to be able to do. These were directly obtained from the research objectives stated in chapter one. The full set of functional requirements are listed in Table~\ref{tab:functional_req} along with the related research objectives.

\begin{table}[H]
\centering
\small
\singlespacing
\renewcommand{\arraystretch}{1.15}
\caption{Functional Requirements}
\label{tab:functional_req}
\begin{tabularx}{\textwidth}{|l|>{\raggedright\arraybackslash}X|c|}
\hline
\textbf{ID} & \textbf{Requirement} & \textbf{Objective} \\
\hline
FR-01 & The system shall allow users to register with email, password, and role selection. & i \\
\hline
FR-02 & The system shall authenticate users using JWT tokens with role-based access control. & i \\
\hline
FR-03 & The system shall allow Sui wallet-based authentication via Ed25519 signature verification. & iii \\
\hline
FR-04 & The system shall capture environmental telemetry data (temperature, soil moisture, soil pH). & i \\
\hline
FR-05 & The system shall generate SHA-256 integrity hashes for all telemetry records. & iii \\
\hline
FR-06 & The system shall create produce batch records linked to telemetry data. & i \\
\hline
FR-07 & The system shall mint produce batch objects on the Sui blockchain. & iii \\
\hline
FR-08 & The system shall track custody transfers of produce batches between stakeholders. & iii \\
\hline
FR-09 & The system shall compute crop yield predictions based on historical telemetry data. & ii \\
\hline
FR-10 & The system shall provide role-specific dashboards for farmers, logistics handlers, and administrators. & i \\
\hline
\end{tabularx}
\end{table}

\clearpage
\subsection{Non-Functional Requirements}

The non-functional requirements specify the qualities of the system and its operational constraints. Table~\ref{tab:nonfunctional_req} shows the non-functional requirements which control system performances, security and usability.

\begin{table}[H]
\centering
\small
\singlespacing
\renewcommand{\arraystretch}{1.15}
\caption{Non-Functional Requirements}
\label{tab:nonfunctional_req}
\begin{tabularx}{\textwidth}{|l|>{\raggedright\arraybackslash}X|}
\hline
\textbf{ID} & \textbf{Requirement} \\
\hline
NFR-01 & All passwords shall be hashed using Argon2id with a minimum cost factor. \\
\hline
NFR-02 & API responses for standard operations shall return within 500 milliseconds. \\
\hline
NFR-03 & The web interface shall be responsive and accessible from modern web browsers. \\
\hline
NFR-04 & Blockchain transactions shall confirm within the Sui Testnet's normal finality window. \\
\hline
NFR-05 & The system shall enforce rate limiting on authentication endpoints (5 attempts per minute). \\
\hline
NFR-06 & Telemetry input values shall be validated within physically plausible ranges. \\
\hline
\end{tabularx}
\end{table}

\section{System Architecture}

The architecture of TerraNode is three-tier, where the presentation layer, the business logic layer and the data persistence layer are separated. This architectural pattern is to achieve modularity so that each layer can be developed, tested, and maintained independently.

\begin{enumerate}[label=(\roman*)]
    \item \textbf{Presentation Layer (Frontend):} Single page application developed using React and TypeScript, to render the user interface and handle client-side state. The front end only communicates with the back end via RESTful API calls, and is connected directly to the Sui blockchain to handle wallet operations via the Mysten dApp Kit.
    \item \textbf{Business Logic Layer (Backend):} Django REST Framework application that handles the API calls, implements business rules, handles user authentication, does analytical prediction calculation, and communicates with the blockchain layer. The API modules exposed to the backend are four domain-specific modules: Authentication, Telemetry, Ledger, and Analytics.
    \item \textbf{Data Persistence Layer:} Three storage units are involved. Structured application data is stored in a primary relational database, PostgreSQL. Redis is an in-memory cache for the most frequently accessed analytics results. The Sui blockchain serves as the unalterable record of produce batch records and their integrity hashes.
\end{enumerate}

The overall system architecture and data flow between the various components is depicted in Figure~\ref{fig:system_arch}.

\begin{figure}[H]
    \centering
    \includegraphics[width=1.0\textwidth]{figures/fig_system_architecture.png}
    \caption{TerraNode System Architecture}
    \label{fig:system_arch}
\end{figure}

\section{Data Flow Design}

\subsection{Context Diagram (Level 0)}

As illustrated in Figure~\ref{fig:dfd_level0}, TerraNode is a single process that communicates with four external processes: the Farmer, the Logistics Handler, the System Administrator and the Sui Blockchain Network.

\begin{figure}[H]
    \centering
    \includegraphics[width=1.0\textwidth]{figures/fig_dfd_level0.png}
    \caption{Context Diagram (DFD Level 0)}
    \label{fig:dfd_level0}
\end{figure}

\subsection{Level 1 Data Flow Diagram}

The Level 1 Data Flow Diagram (Figure~\ref{fig:dfd_level1}) breaks down TerraNode into four main processes:

\begin{enumerate}[label=(\roman*)]
    \item \textbf{Process 1.0 --- Authentication Module:} Responsible for registering users, logging in using credentials, signing in with wallet and managing sessions using JWT tokens.
    \item \textbf{Process 2.0 --- Telemetry Module:} Reads data from environmental sensors, checks the range of the input data, produces the integrity hash of the input data using the SHA-256 algorithm, and saves the data in the database.
    \item \textbf{Process 3.0 --- Ledger Module:} Responsible for produce batch creation, working with the Sui blockchain to mint and transfer the produce to the chain, and coordinating with the Sui blockchain to synchronize the local database with the chain.
    \item \textbf{Process 4.0 --- Analytics Module:} This module fetches historical telemetry data, uses the weighted moving average prediction algorithm, stores the predictions in Redis cache, and returns the yield forecast to the user.
\end{enumerate}

\begin{figure}[H]
    \centering
    \includegraphics[width=1.0\textwidth]{figures/fig_dfd_level1.png}
    \caption{Data Flow Diagram (Level 1)}
    \label{fig:dfd_level1}
\end{figure}

\section{Entity-Relationship Model}

The TerraNode database is composed of three main entities and has well-defined attributes and relationships. The Entity-Relationship (ER) diagram is given in Figure~\ref{fig:er_diagram}.

\begin{figure}[H]
    \centering
    \includegraphics[width=1.0\textwidth]{figures/fig_er_diagram.png}
    \caption{Entity-Relationship Diagram}
    \label{fig:er_diagram}
\end{figure}

All users of the system are represented by \textbf{CustomUser} which is based on Django's AbstractUser model. A UUID primary key identifies each user and they are authenticated by their email address. The three types of users are: Farmers, Logistics Handlers and Administrators, which are differentiated using the role field. The wallet-based authentication uses a wallet address provided in an optional \texttt{sui\_public\_key} field which stores the user's Sui blockchain wallet address.

EnvironmentalTelemetry: Data for single environmental readings recorded by farmers. The record contains temperature in Celsius, soil moisture (%) and soil pH. The \texttt{payload\_sha256} field contains the sha256 hash of the canonical representation of the record so that any changes made after it has been created can be detected. There is a foreign key relationship created between each telemetry record and the farmer who submitted it.

\textbf{ProduceBatch} is a batch of agricultural produce that passes through the supply chain. Each batch is associated to the farmer who made it, optionally associated to a telemetry record which describes the results of the environment under which the produce had been grown and has a \texttt{current\_custodian} field which tracks the current physical owner of the batch. The batch lifecycle is controlled by the status field, which can have four different statuses: PENDING, MINTED, IN\_TRANSIT, DELIVERED. The Sui network returns identifiers on successful on-chain operations, with two specific fields: \texttt{sui\_object\_id} and \texttt{sui\_tx\_digest} for blockchain operations.

\section{Use Case Modelling}

The use case diagram in figure~\ref{fig:use_case} shows the main interactions among the three types of actors and the TerraNode system.

\begin{figure}[H]
    \centering
    \includegraphics[width=1.0\textwidth]{figures/fig_use_case.png}
    \caption{Use Case Diagram}
    \label{fig:use_case}
\end{figure}

\subsection{Actor Descriptions}

The \textbf{Farmer} actor can register an account, login (with credentials or wallet), upload telemetry readings, add produce batches, mint batches on the Sui blockchain, check the yield predictions and adjust their settings.

The \textbf{Logistics Handler} actor can log-in, get batches of produce through on-chain transfers, update the shipment statuses and see where and how batches are being taken care of.

The \textbf{System Administrator} actor has access to all data in regards to the system, can control user accounts, check the audit log and monitor the system infrastructure.

\section{Smart Contract Design}

The TerraNode smart contract has been deployed on the Sui Testnet and is written in the Move programming language. It outlines the definition of batches of produce on the chain and the actions that can be taken with them.

\subsection{ProduceBatch Object Structure}

The \texttt{ProduceBatch} struct is defined as a Sui object that can be stored, owned, and transferred on-chain, and has the \texttt{key} and \texttt{store} abilities. It has the following fields:

\begin{enumerate}[label=(\roman*)]
    \item \texttt{id} (UID): Unique object identifier provided by the Sui runtime.
    \item \texttt{crop\_type} (String): The type of a crop that is included in the batch.
    \item \texttt{weight\_kg} (u64): The weight of the batch in kg.
    \item \texttt{origin\_farmer\_address} (address): The Sui address of the owner of the batch.
    \item \texttt{current\_custodian\_address} (address): Sui address of the current custodian of the batch.
    \item \texttt{data\_integrity\_hash} (vector$<$u8$>$): The hash of the data, computed using the SHA-256 algorithm, that guarantees the data source that the data in the on-chain record is the same as the data they originally sent.
\end{enumerate}

\subsection{Entry Functions}

The smart contract has two entry functions:

The function \texttt{mint\_batch} creates a new \texttt{ProduceBatch} object, sets the origin farmer and the first custodian as the transaction sender and emits a \texttt{BatchMinted} event. The batch object is then sent to the sender's address.

\textbf{transfer\_custody:} Takes a modifiable reference to an existing \texttt{ProduceBatch} object and an address for the new custodian. Before updating the \texttt{current\_custodian\_address} field and emitting a \texttt{CustodyTransferred} event, the function checks that the caller is the current custodian.
\subsection{Event Emissions}

Two types of events are emitted by the smart contract for off-chain indexing: \texttt{BatchMinted} which contains the batch identifier and the farmer's address and the crop type and \texttt{CustodyTransferred} which contains the batch identifier, previous custodian address, and new custodian address. The Sui JSON-RPC API allows these events to be queried, allowing the TerraNode backend to stay in sync with the state on-chain.

\begin{figure}[H]
    \centering
    \includegraphics[width=1.0\textwidth]{figures/fig_batch_lifecycle.png}
    \caption{Batch Lifecycle Sequence Diagram}
    \label{fig:batch_lifecycle}
\end{figure}

\section{Predictive Analytics Engine}

\subsection{Feature Vector Composition}

The analytics engine takes three environmental variables extracted from the records of the EnvironmentalTelemetry, which are temperature in degrees Celsius ($T$), soil moisture as a percentage ($M$), and soil pH ($P$), as its input and feature vector. If the same farmer has multiple telemetry records, the engine fetches the last 90 records in descending order by telemetry time.

\subsection{Weighted Moving Average Model}

The current implementation was based on a simplified weighted moving average model to estimate the yield of the crop. For $n$ telemetry records, the model first computes the arithmetic means of the three environmental features:

\begin{equation}
\bar{T} = \frac{1}{n} \sum_{i=1}^{n} T_i, \quad
\bar{M} = \frac{1}{n} \sum_{i=1}^{n} M_i, \quad
\bar{P} = \frac{1}{n} \sum_{i=1}^{n} P_i
\label{eq:averages}
\end{equation}

The predicted yield $Y$ in metric tons is then computed as:

\begin{equation}
Y = \max\left(0, \; 50.0 + 0.5 \cdot \bar{M} - 2.0 \cdot |22.0 - \bar{T}|\right)
\label{eq:yield}
\end{equation}

50.0 is a nominal base yield. The moisture coefficient (0.5) indicates that crop productivity increases in a linear manner as the amount of soil moisture available increases. The temperature penalty term ($2.0 \cdot |22.0 - \bar{T}|$) is a negative term that represents the deviation from optimal growing temperature (22 degrees Celsius). The $\max(0, \ldots)$ function makes sure to always have a positive predicted yield.

\subsection{Confidence Score}

The confidence score $C$ reflects the data density available for the prediction:

\begin{equation}
C = \min\left(0.95, \; \frac{n}{90}\right)
\label{eq:confidence}
\end{equation}

The confidence score is monotonically increasing with the amount of data, with a maximum value of 0.95 reflecting the fact that the more historical data that has been observed, the more confident one can be in the prediction. The engine will return an error if only a less number of data points are provided than it requires, which is five.

\subsection{Caching Strategy}

The prediction results are persisted in Redis with TTL of 10 minutes (600 seconds). The cache key is created in the following way: \texttt{analytics:predict:}\allowbreak\texttt{\{farmer\_id\}}. Upon receiving a prediction request, the engine will check the cache first. If there is a valid cached result, it is returned immediately. Otherwise, it is computed, cached in memory and then returned to the client.

\section{Security Architecture}

\subsection{Authentication}

Two levels of authentication are provided at TerraNode. The main approach is by email and password, which is handled by the Django REST Framework (DRF) library Simple JWT. If authentication is successful, a pair of tokens are issued: an access token (default lifetime: 15 minutes) and a refresh token (default lifetime: 7 days). Each use of a refresh token, a new pair is generated and the old refresh token is blacklisted, which is called token rotation.

\subsection{Sui Wallet Signature Verification}
\label{sec:wallet_verify}

Users will be able to log in with their Sui blockchain wallet through the secondary authentication method. The verification process is as follows:

\begin{enumerate}[label=(\roman*)]
    \item The client builds a Login message with the form: \\ ``\texttt{Login to TerraNode with \{sui\_public\_key\}}''.
    \item This message is signed by the user with his/her private key in the Sui wallet.
    \item The message, signature and public key are sent to the backend.
    \item The backend is able to reconstruct the signed payload by prepending the Sui personal message intent prefix (\texttt{[3, 0, 0]}) and encoding the message length with ULEB128.
    \item The payload is reconstructed and then hashed with BLAKE2b 32-byte digest.
    \item The payload is then hashed and a signature is verified with the provided public key.
    \item If the verification is successful, and the public key is matched with a registered user, JWT tokens are issued.
\end{enumerate}

\subsection{Data Integrity Hashing}

Each telemetry record has a SHA256 hash, which is derived from a common, standardized JSON representation of its essential data. With the canonical form, the keys are sorted, there is no whitespace, and the formatting of the floating point is fixed to guarantee deterministic output. This is a hash which is stored with the record and which can be independently calculated at any time to confirm that the record is not changed since it was created. This hash is also included in the Sui blockchain as part of the respective ProduceBatch object, establishing an on-chain link to the data source in the real world.

\subsection{Role-Based Access Control}

The backend does role based access control (RBAC) using custom permission classes. There are four levels of permission:

\begin{enumerate}[label=(\roman*)]
    \item \textbf{IsFarmer:} Provides access to endpoints that generate telemetry records, batch produce and monitor yield predictions.
    \item \textbf{IsLogistics:} Provides batch transfer and shipment management endpoints.
    \item \textbf{IsAdmin:} Gives access to administrative operations such as user management, and viewing of the audit log.
    \item \textbf{IsFarmerOrAdmin:} Provides shared access to endpoints for farmers and admin, such as viewing telemetry history.
\end{enumerate}

Figure~\ref{fig:auth_flow} illustrates the authentication flow for both credential-based and wallet-based login.

\begin{figure}[H]
    \centering
    \includegraphics[width=1.0\textwidth]{figures/fig_auth_flow.png}
    \caption{Authentication Flow Diagram}
    \label{fig:auth_flow}
\end{figure}

\section{Technology Stack}
\label{sec:tech_stack}

As shown in Table~\ref{tab:tech_stack}, TerraNode is developed with the following technologies.

\begin{table}[H]
\centering
\small
\singlespacing
\renewcommand{\arraystretch}{1.15}
\caption{Technology Stack}
\label{tab:tech_stack}
\begin{tabularx}{\textwidth}{|l|l|>{\raggedright\arraybackslash}X|}
\hline
\textbf{Layer} & \textbf{Technology} & \textbf{Purpose} \\
\hline
Frontend & React 19 & Component-based user interface library \\
\hline
Frontend & TypeScript 6.0 & Statically typed JavaScript superset \\
\hline
Frontend & Vite 8.1 & Build tool and development server \\
\hline
Frontend & Recharts 3.9 & Data visualization (charts and graphs) \\
\hline
Frontend & Mysten dApp Kit & Sui wallet integration and transaction signing \\
\hline
Backend & Django 5.1 & Python web framework \\
\hline
Backend & Django REST Framework 3.15 & RESTful API toolkit \\
\hline
Backend & Simple JWT 5.3 & JSON Web Token authentication \\
\hline
Backend & Celery & Asynchronous task queue \\
\hline
Database & PostgreSQL & Primary relational database \\
\hline
Cache & Redis 5.0 & In-memory data store for caching \\
\hline
Blockchain & Sui Testnet & Delegated PoS public blockchain \\
\hline
Smart Contract & Move & Smart contract programming language \\
\hline
Security & Argon2id & Password hashing algorithm \\
\hline
Security & PyNaCl & Ed25519 signature verification \\
\hline
\end{tabularx}
\end{table}
\section{Development Tools}

The system has been developed using the Visual Studio Code as the main integrated development environment. Version control was maintained using Git throughout the entire project. The front-end was installed and managed using npm (package manager) and the back-end was installed and managed using Pipenv for python dependency isolation. Testing was done with Django's built-in test framework with pytest added and Django factory-boy and faker were used to generate test data.

\section{Chapter Summary}

In this chapter, the methodology used for the design and development of TerraNode was introduced. This approach to development, called Design Science Research, was explained, along with a description of functional and non-functional requirements. Explanation of the three-tier system architecture, and first level (context) and second level (decomposition) DFD's were provided. The entity-relationship model determined the main data structures, and the use case diagram described the interactions among the three types of actors and the system. The smart contract design was documented including the object structure, entry functions and event emissions. The predictive analytics engine has been mathematically formulated and formally described with the mathematical formulations for yield prediction and confidence scoring. Lastly, security architecture was provided, such as JWT authentication, Sui wallet verification, SHA-256 data integrity hashing, and role-based access control. This design implementation is described in the next chapter.